\documentclass[10pt,letterpaper]{article}

\usepackage[
  letterpaper,
  top=0.7in,
  bottom=0.7in,
  left=0.8in,
  right=0.8in
]{geometry}
\usepackage[T1]{fontenc}
\usepackage{mathptmx}
\usepackage{amsmath}
\usepackage{booktabs}
\usepackage{array}
\usepackage{cite}
\usepackage{url}
\usepackage{float}

\setlength{\parindent}{0.14in}
\setlength{\parskip}{0pt}
\pagestyle{plain}
\renewcommand{\thesection}{\Roman{section}}
\renewcommand{\thesubsection}{\Alph{subsection}.}
\renewcommand{\thetable}{\Roman{table}}
\renewcommand{\refname}{REFERENCES}

\makeatletter
\renewcommand\section{\@startsection{section}{1}{\z@}%
  {1.3ex plus 1ex minus .2ex}%
  {0.8ex plus .2ex}%
  {\centering\normalfont\small\scshape}}
\renewcommand\subsection{\@startsection{subsection}{2}{\z@}%
  {1.0ex plus .8ex minus .2ex}%
  {0.6ex plus .2ex}%
  {\normalfont\itshape}}
\renewcommand{\fnum@table}{\textsc{Table \thetable}}
\makeatother

\begin{document}

\begin{center}
{\LARGE Lightweight RGB Hand Pose Tracking for Edge HCI\par}
\vspace{0.6em}
{\large Brandon MacGillivray\par}
\vspace{0.2em}
{\normalsize School of Electronics, Electrical Engineering and Computer Science\par}
{\normalsize Queen's University Belfast, Belfast, UK\par}
{\normalsize \texttt{bmacgillivray01@qub.ac.uk}\par}
\end{center}

\vspace{0.6em}
\noindent\textbf{\textit{Abstract}}---

\vspace{0.6em}
\noindent\textbf{\textit{Keywords}}---

\section{INTRODUCTION}

- why is it worthwhile to work on hand pose estimation?

\section{LITERATURE REVIEW}

- what are knowledge gaps based on existing literatre 

\section{TECHNICAL APPROACH}

\subsection{Project Scope}
- what questions am i trying to answer? \\
- Can a reduced 10-keypoint model maintain sufficient hand-pose accuracy for edge HCI while improving runtime efficiency? HCI does not need 21 keypoints to be calcualted as not all are used for gesture calculation, thus wasted computation.\\
- Does dual-branch fusion outperform heatmap-only and coordinate-only prediction? from the outset, it seems like DRHand is wasting computation having two branches, is it wasted?\\
- To what extent can transfer learning improve cross-dataset hand-pose performance between RHD and COCO-style hand data? caused by the lack of real data for training this model. \\
- Is the proposed pipeline competitive enough in accuracy and latency for practical embedded deployment compared with a baseline such as MediaPipe? is it useful?\\

\subsection{DRHand-Inspired Model}

- in brief, how does DRHand work? 
- network structure 
- branches
- post processing
- loss


\subsection{Relation to the Original DRHand Paper}

- how is this implementation different to orginal paper?

\subsection{Experimental Setup}

- dataset detail?
- implementation details?  

\subsection{Evaluation Metrics}

- how are the models evalulated?
- what is the math?
- what doe the metrics mean?

\section{RESULTS AND ANALYSIS}

\subsection{Experimental Results}

\noindent\textbf{Point:} The reduced 10-keypoint model preserves most of the useful pose-estimation performance of the 21-keypoint baseline on the fair shared-subset comparison.
\textbf{Evidence:} In Table~\ref{tab:main_accuracy}, \texttt{B0} shared-10 gives SSE 0.1511, EPE 0.0764, and PCK 0.9101, while \texttt{B1} native-10 gives SSE 0.1574, EPE 0.0706, and PCK 0.9099.
\textbf{Explanation:} The two models are therefore very close: \texttt{B0} is slightly better on SSE, \texttt{B1} is slightly better on EPE, and PCK is effectively unchanged.
\textbf{Link:} This supports the research question that a reduced landmark layout can maintain sufficient accuracy for the intended edge HCI setting.

\medskip
\noindent\textbf{Point:} The runtime advantage of the 10-keypoint model appears to arise mainly from fusion post-processing rather than from the core branch inference paths.
\textbf{Evidence:} In Table~\ref{tab:latency}, the forward-time gap between \texttt{B0} and \texttt{B1} is large in fusion mode (27.43 vs 24.59 ms) but small in heatmap mode (24.28 vs 24.00 ms) and coordinate mode (24.08 vs 23.91 ms).
\textbf{Explanation:} Since the single-branch modes show little latency difference, the main extra cost in the 21-joint case is likely introduced by the heavier fusion stage.
\textbf{Link:} This refines the efficiency result by showing that the 10-keypoint speedup is concentrated primarily in fusion mode.

\medskip
\noindent\textbf{Point:} Fusion remains the strongest overall operating mode for the proposed architecture.
\textbf{Evidence:} In Table~\ref{tab:branch_ablation}, fusion gives the best overall balance of SSE and EPE for both \texttt{B0} and \texttt{B1}, while heatmap-only and coordinate-only outputs are less consistent.
\textbf{Explanation:} This indicates that the heatmap and coordinate branches provide complementary information rather than duplicating the same function.
\textbf{Link:} Therefore, the results support the research question that the dual-branch design is justified.

\medskip
\noindent\textbf{Point:} Transfer learning is effective when used for target-domain fine-tuning, but not as a zero-shot cross-dataset solution.
\textbf{Evidence:} In Table~\ref{tab:transfer_learning}, \texttt{C\_ONLY} performs poorly on RHD, but \texttt{C\_TO\_R} improves strongly after fine-tuning; similarly, \texttt{R\_ONLY} performs poorly on COCO, but \texttt{R\_TO\_C} improves substantially after adaptation.
\textbf{Explanation:} This shows that source-domain training learns useful reusable features, but these are not sufficient on their own without target-domain adjustment.
\textbf{Link:} The transfer results therefore answer the research question by showing that transfer is useful as initialization for fine-tuning rather than as a direct replacement for target-domain training.

\medskip
\noindent\textbf{Point:} Transfer performance is asymmetric across datasets.
\textbf{Evidence:} Table~\ref{tab:transfer_learning} shows that \texttt{R\_TO\_C} approaches COCO scratch performance much more closely than \texttt{C\_TO\_R} approaches RHD scratch performance.
\textbf{Explanation:} This suggests that transfer difficulty depends on the source-target pairing and that the two datasets are not equally compatible.
\textbf{Link:} Accordingly, the transfer study shows not only that fine-tuning helps, but also that some directions of adaptation are easier than others.

% this seems sus. remove the handiness results as that hinders mediapipe alot. also I wonder if the --right flag is a problem. might have to run again using --auto for mp to be fair
\medskip
\noindent\textbf{Point:} The learned DRHand-style models outperform the MediaPipe baseline on RHD.
\textbf{Evidence:} In Table~\ref{tab:mediapipe_baseline}, MediaPipe gives much worse SSE and PCK than the learned models in both native-21 and shared-10 comparisons, while ignoring handedness changes its detection success markedly.
\textbf{Explanation:} This indicates that the learned pipeline is substantially stronger on the benchmark dataset and less dependent on the particular handedness filtering choice used in the baseline.
\textbf{Link:} Therefore, the proposed system can be considered competitive and practically useful relative to a relevant external baseline.

% no it's not, as the accuracy is not high enough, this could be seen in the qualitative results if I had any. must wait on that to answer this question. 
% only reasonable based on latency results
\medskip
\noindent\textbf{Point:} The proposed pipeline is viable for practical embedded deployment.
\textbf{Evidence:} The Jetson results in Table~\ref{tab:latency} show throughput of roughly 27 to 31 images per second across the main operating modes.
\textbf{Explanation:} These rates are close to real-time performance, which means the system is not only accurate in offline evaluation but also feasible for edge-device use.
\textbf{Link:} This directly supports the research question concerning practical embedded deployment.

% how could there be such a low oracle score yet the aggregated results still be good? seems sus.
\medskip
\noindent\textbf{Point:} The current fusion heuristic remains a clear weakness of the system.
\textbf{Evidence:} In Table~\ref{tab:fusion_diagnostics}, the oracle-match rate is only around 0.48 to 0.53.
\textbf{Explanation:} This means the current rule often fails to choose the lower-error branch, even though fusion still improves aggregate performance.
\textbf{Link:} Consequently, the experiments identify fusion design as the clearest target for future refinement.

\subsection{Qualitative Results}

- test on small set if real images manually annotated by myself\\
- 21 vs 10 on RHD, coco, RHD to coco, coco to RHD, mp

\subsection{Tables}

\begin{table}[H]
\centering
\footnotesize
\caption{Main quantitative comparison between the 21-keypoint baseline and the reduced 10-keypoint model on RHD. Lower SSE and EPE are better; higher PCK@0.2 is better. Values are mean $\pm$ std over five seeds.}
\label{tab:main_accuracy}
\begin{tabular}{llcccc}
\toprule
Model & Eval set & Train joints & SSE $\downarrow$ & EPE $\downarrow$ & PCK@0.2 $\uparrow$ \\
\midrule
B0 baseline & native-21 & 21 & 0.3098 $\pm$ 0.0236 & 0.0779 $\pm$ 0.0054 & 0.9112 $\pm$ 0.0071 \\
B0 baseline & shared-10 & 21 & 0.1511 $\pm$ 0.0120 & 0.0764 $\pm$ 0.0039 & 0.9101 $\pm$ 0.0071 \\
B1 reduced & native-10 & 10 & 0.1574 $\pm$ 0.0226 & 0.0706 $\pm$ 0.0054 & 0.9099 $\pm$ 0.0046 \\
\bottomrule
\end{tabular}
\end{table}

\begin{table}[H]
\centering
\footnotesize
\caption{Branch ablation for the dual-regression design. Fusion is compared against the heatmap-only and coordinate-only outputs using the same trained checkpoints. Values are mean $\pm$ std over five seeds.}
\label{tab:branch_ablation}
\begin{tabular}{llcccc}
\toprule
Model & Eval set & Mode & SSE $\downarrow$ & EPE $\downarrow$ & PCK@0.2 $\uparrow$ \\
\midrule
B0 & native-21 & fusion & 0.3098 $\pm$ 0.0236 & 0.0779 $\pm$ 0.0054 & 0.9112 $\pm$ 0.0071 \\
B0 & native-21 & heatmap & 0.3139 $\pm$ 0.0684 & 0.0989 $\pm$ 0.0335 & 0.9374 $\pm$ 0.0066 \\
B0 & native-21 & coord & 0.3134 $\pm$ 0.0229 & 0.0788 $\pm$ 0.0062 & 0.9113 $\pm$ 0.0071 \\
B1 & native-10 & fusion & 0.1574 $\pm$ 0.0226 & 0.0706 $\pm$ 0.0054 & 0.9099 $\pm$ 0.0046 \\
B1 & native-10 & heatmap & 0.2213 $\pm$ 0.0880 & 0.0991 $\pm$ 0.0288 & 0.9022 $\pm$ 0.0393 \\
B1 & native-10 & coord & 0.1599 $\pm$ 0.0218 & 0.0715 $\pm$ 0.0070 & 0.9102 $\pm$ 0.0048 \\
\bottomrule
\end{tabular}
\end{table}

\begin{table}[H]
\centering
\footnotesize
\caption{Transfer-learning results for the four training sequences evaluated on both RHD and the COCO-style hand dataset. Values are mean $\pm$ std over five seeds, and fusion-mode evaluation is shown because it is the main operating mode in the study.}
\label{tab:transfer_learning}
\begin{tabular}{llcccccc}
\toprule
Training sequence & Exp. & RHD SSE & RHD EPE & RHD PCK & COCO SSE & COCO EPE & COCO PCK \\
\midrule
RHD scratch & \texttt{R\_ONLY} & 0.3091 $\pm$ 0.0269 & 0.0764 $\pm$ 0.0030 & 0.9143 $\pm$ 0.0097 & 3.7488 $\pm$ 0.7985 & 0.3365 $\pm$ 0.0075 & 0.3719 $\pm$ 0.0434 \\
COCO scratch & \texttt{C\_ONLY} & 1.4194 $\pm$ 0.2619 & 0.1609 $\pm$ 0.0207 & 0.3397 $\pm$ 0.0718 & 0.0764 $\pm$ 0.0016 & 0.0612 $\pm$ 0.0049 & 0.9834 $\pm$ 0.0007 \\
COCO $\rightarrow$ RHD & \texttt{C\_TO\_R} & 0.4504 $\pm$ 0.0162 & 0.1020 $\pm$ 0.0039 & 0.8190 $\pm$ 0.0114 & 1.6426 $\pm$ 0.2293 & 0.2678 $\pm$ 0.0114 & 0.5320 $\pm$ 0.0164 \\
RHD $\rightarrow$ COCO & \texttt{R\_TO\_C} & 0.5699 $\pm$ 0.0592 & 0.1276 $\pm$ 0.0149 & 0.7690 $\pm$ 0.0221 & 0.0980 $\pm$ 0.0064 & 0.0840 $\pm$ 0.0018 & 0.9837 $\pm$ 0.0007 \\
\bottomrule
\end{tabular}
\end{table}

\begin{table}[H]
\centering
\footnotesize
\caption{Comparison between the MediaPipe Hand Landmarker baseline and the learned DRHand variants on RHD. The MediaPipe rows are single evaluation runs, whereas the DRHand rows are mean $\pm$ std over five seeds. Detection success is shown only for the MediaPipe runs because the learned-model evaluations do not use the same detection stage.}
\label{tab:mediapipe_baseline}
\begin{tabular}{llcccc}
\toprule
Method & Eval set & SSE & EPE & PCK@0.2 & Detect. succ. \\
\midrule
MediaPipe & native-21 & 2.3022 & 0.0830 & 0.6452 & 0.4879 \\
MediaPipe (ignore handedness) & native-21 & 1.4002 & 0.0904 & 0.6726 & 0.8933 \\
MediaPipe & shared-10 & 1.1202 & 0.1209 & 0.6407 & 0.4879 \\
\texttt{B0} fusion & native-21 & 0.3098 $\pm$ 0.0236 & 0.0779 $\pm$ 0.0054 & 0.9112 $\pm$ 0.0071 & n/a \\
\texttt{B0} fusion & shared-10 & 0.1511 $\pm$ 0.0120 & 0.0764 $\pm$ 0.0039 & 0.9101 $\pm$ 0.0071 & n/a \\
\texttt{B1} fusion & native-10 & 0.1574 $\pm$ 0.0226 & 0.0706 $\pm$ 0.0054 & 0.9099 $\pm$ 0.0046 & n/a \\
\bottomrule
\end{tabular}
\end{table}

\begin{table}[H]
\centering
\footnotesize
\caption{Edge-device latency summary for the main models. Forward time measures the prediction stage only, whereas end-to-end time includes image read, preprocessing, host-to-device transfer, prediction, and JSON write. Values are mean $\pm$ std over five seeds.}
\label{tab:latency}
\begin{tabular}{llccc}
\toprule
Model & Mode & Forward ms/img $\downarrow$ & End-to-end ms/img $\downarrow$ & Images/s $\uparrow$ \\
\midrule
B0 & fusion & 27.43 $\pm$ 0.45 & 36.81 $\pm$ 1.33 & 27.20 $\pm$ 0.97 \\
B0 & heatmap & 24.28 $\pm$ 0.35 & 33.19 $\pm$ 0.56 & 30.13 $\pm$ 0.50 \\
B0 & coord & 24.08 $\pm$ 0.44 & 33.17 $\pm$ 0.86 & 30.16 $\pm$ 0.77 \\
B1 & fusion & 24.59 $\pm$ 0.61 & 33.31 $\pm$ 0.71 & 30.03 $\pm$ 0.64 \\
B1 & heatmap & 24.00 $\pm$ 0.34 & 32.64 $\pm$ 0.46 & 30.64 $\pm$ 0.43 \\
B1 & coord & 23.91 $\pm$ 0.29 & 32.50 $\pm$ 0.34 & 30.77 $\pm$ 0.32 \\
\bottomrule
\end{tabular}
\end{table}

\begin{table}[H]
\centering
\footnotesize
\caption{Detailed latency breakdown for the main fused models. These metrics are available from the benchmark pipeline and help localize where runtime is spent. Values are mean $\pm$ std over five seeds.}
\label{tab:latency_breakdown}
\begin{tabular}{lcccccc}
\toprule
Model & Read ms & Prep ms & H2D ms & Forward ms & Write ms & E2E ms \\
\midrule
B0 fusion & 4.89 $\pm$ 0.64 & 2.80 $\pm$ 0.13 & 0.32 $\pm$ 0.02 & 27.43 $\pm$ 0.45 & 0.67 $\pm$ 0.09 & 36.81 $\pm$ 1.33 \\
B1 fusion & 4.49 $\pm$ 0.08 & 2.74 $\pm$ 0.05 & 0.31 $\pm$ 0.01 & 24.59 $\pm$ 0.61 & 0.52 $\pm$ 0.02 & 33.31 $\pm$ 0.71 \\
\bottomrule
\end{tabular}
\end{table}

\begin{table}[H]
\centering
\footnotesize
\caption{Fusion diagnostics for the dual-regression post-processing rule. The oracle-match rate measures how often the final fusion choice matches the lower-error branch for visible joints. Values are mean $\pm$ std over five seeds. The table includes all currently available diagnostic exports, including a shared-10 row for \texttt{B0}.}
\label{tab:fusion_diagnostics}
\begin{tabular}{llccccc}
\toprule
Model & Eval set & HM sel. & Coord sel. & $\alpha$ mean & Disagr. mean & Oracle match \\
\midrule
B0 & native-21 & 0.426 $\pm$ 0.025 & 0.574 $\pm$ 0.025 & 0.065 $\pm$ 0.007 & 0.163 $\pm$ 0.022 & 0.481 $\pm$ 0.012 \\
B0 & shared-10 & 0.407 $\pm$ 0.024 & 0.593 $\pm$ 0.024 & 0.065 $\pm$ 0.007 & 0.162 $\pm$ 0.018 & 0.474 $\pm$ 0.011 \\
B1 & native-10 & 0.501 $\pm$ 0.053 & 0.499 $\pm$ 0.053 & 0.084 $\pm$ 0.003 & 0.171 $\pm$ 0.018 & 0.529 $\pm$ 0.037 \\
\bottomrule
\end{tabular}
\end{table}

\begin{table}[H]
\centering
\footnotesize
\caption{Hyperparameter ablation relative to the baseline \texttt{B0} configuration, using native-21 fused accuracy on RHD and fused end-to-end benchmark latency. Values are mean $\pm$ std over five seeds.}
\label{tab:hyper_ablation}
\begin{tabular}{llcccc}
\toprule
Exp & Change vs. B0 & SSE $\downarrow$ & EPE $\downarrow$ & PCK@0.2 $\uparrow$ & E2E ms $\downarrow$ \\
\midrule
B0 & baseline & 0.3098 $\pm$ 0.0236 & 0.0779 $\pm$ 0.0054 & 0.9112 $\pm$ 0.0071 & 36.81 $\pm$ 1.33 \\
L1 & $\lambda_{hm}=1.5$, $\lambda_{coord}=0.5$ & 0.3583 $\pm$ 0.0534 & 0.0725 $\pm$ 0.0096 & 0.9081 $\pm$ 0.0092 & 35.77 $\pm$ 0.45 \\
L2 & $\lambda_{hm}=0.5$, $\lambda_{coord}=1.5$ & 0.3073 $\pm$ 0.0260 & 0.0812 $\pm$ 0.0022 & 0.9115 $\pm$ 0.0095 & 35.75 $\pm$ 0.49 \\
S1 & $\sigma=1.5$ & 0.3252 $\pm$ 0.0381 & 0.0707 $\pm$ 0.0055 & 0.9116 $\pm$ 0.0086 & 35.72 $\pm$ 0.25 \\
S2 & $\sigma=3.0$ & 0.3137 $\pm$ 0.0488 & 0.0833 $\pm$ 0.0058 & 0.9177 $\pm$ 0.0081 & 36.48 $\pm$ 0.53 \\
W1 & $w=5.0$ & 0.3213 $\pm$ 0.0384 & 0.0737 $\pm$ 0.0058 & 0.9150 $\pm$ 0.0099 & 36.21 $\pm$ 0.39 \\
W2 & $w=15.0$ & 0.2979 $\pm$ 0.0441 & 0.0750 $\pm$ 0.0031 & 0.9140 $\pm$ 0.0124 & 35.91 $\pm$ 0.43 \\
E1 & $\epsilon=1.0$ & 0.2994 $\pm$ 0.0591 & 0.0695 $\pm$ 0.0011 & 0.9175 $\pm$ 0.0113 & 35.84 $\pm$ 0.23 \\
E2 & $\epsilon=4.0$ & 0.3106 $\pm$ 0.0106 & 0.0873 $\pm$ 0.0056 & 0.9110 $\pm$ 0.0067 & 36.13 $\pm$ 0.52 \\
\bottomrule
\end{tabular}
\end{table}

\subsection{Qualitative Results}

- on the same set of real images annotated by me. how did various models perform?
- 21, 10, coco, RHD->coco, coco->RHD, mp

\section{CONCLUSION AND FUTURE WORK}

\clearpage
\begin{thebibliography}{00}

\bibitem{vuletic}
T. Vuletic, A. Duffy, L. Hay, C. McTeague, and G. Campbell, ``Systematic literature review of hand gestures used in human-computer interaction interfaces,'' \emph{Int. J. Human-Computer Studies}, vol. 129, pp. 47--65, 2019.

\bibitem{wachs}
J. P. Wachs, M. Kolsch, H. Stern, and Y. Edan, ``Vision-based hand-gesture applications,'' \emph{Communications of the ACM}, vol. 54, no. 2, pp. 60--71, 2011.

\bibitem{li_vr}
Y. Li, J. Huang, F. Tian, H.-A. Wang, and G.-Z. Dai, ``Gesture interaction in virtual reality,'' \emph{Virtual Reality \& Intelligent Hardware}, vol. 1, no. 1, pp. 84--112, 2019.

\bibitem{dualreg}
D. Wei, S. An, X. Zhang, J. Tian, K. A. Tsintotas, A. Gasteratos, and H. Zhu, ``Dual regression for efficient hand pose estimation,'' in \emph{Proc. IEEE Int. Conf. Robotics and Automation}, 2022.

\bibitem{mediapipe}
F. Zhang, V. Bazarevsky, A. Vakunov, A. Tkachenka, G. Sung, C.-L. Chang, and M. Grundmann, ``MediaPipe hands: On-device real-time hand tracking,'' \emph{arXiv preprint arXiv:2006.10214}, 2020.

\bibitem{srhandnet}
Y. Wang, B. Zhang, and C. Peng, ``SRHandNet: Real-time 2D hand pose estimation with simultaneous region localization,'' \emph{IEEE Trans. Image Processing}, vol. 29, pp. 2977--2986, 2019.

\bibitem{wingloss}
Z.-H. Feng, J. Kittler, M. Awais, P. Huber, and X.-J. Wu, ``Wing loss for robust facial landmark localisation with convolutional neural networks,'' in \emph{Proc. IEEE Conf. Computer Vision and Pattern Recognition}, 2018, pp. 2235--2245.

\bibitem{fasthand}
S. An, X. Zhang, D. Wei, H. Zhu, J. Yang, and K. A. Tsintotas, ``FastHand: Fast monocular hand pose estimation on embedded systems,'' \emph{Journal of Systems Architecture}, vol. 122, p. 102361, 2022.

\bibitem{attentionlite}
N. Santavas, I. Kansizoglou, L. Bampis, E. Karakasis, and A. Gasteratos, ``Attention! A lightweight 2D hand pose estimation approach,'' \emph{IEEE Sensors Journal}, vol. 21, no. 10, pp. 11488--11501, 2021.

\bibitem{stereojetson}
E. Williams-Linera and J. M. Ramirez-Cortes, ``Stereo vision system based on the NVIDIA Jetson Nano for real-time evaluation of yoga poses,'' in \emph{Proc. IEEE}, 2024, pp. 1--7.

\bibitem{blazepose}
V. Bazarevsky, I. Grishchenko, K. Raveendran, T. Zhu, F. Zhang, and M. Grundmann, ``BlazePose: On-device real-time body pose tracking,'' \emph{arXiv preprint arXiv:2006.10204}, 2020.

\end{thebibliography}

\end{document}
